\documentclass[12pt, letterpaper]{article}

% ─── Paquetes ───────────────────────────────────────────────────────────────
\usepackage[spanish]{babel}
\usepackage[utf8]{inputenc}
\usepackage[T1]{fontenc}
\usepackage{geometry}
\usepackage{graphicx}
\usepackage{xcolor}
\usepackage{tikz}
\usepackage{mdframed}
\usepackage{enumitem}
\usepackage{fontawesome5}
\usepackage{hyperref}
\usepackage{fancyhdr}
\usepackage{titlesec}
\usepackage{tcolorbox}
\usepackage{booktabs}
\usepackage{array}
\usepackage{parskip}
\usepackage{lmodern}
\usepackage{microtype}
\usepackage{float}
\usepackage{caption}

\tcbuselibrary{skins, breakable, listings}
\usetikzlibrary{shadows}

% ─── Geometría ──────────────────────────────────────────────────────────────
\geometry{
  top=2.5cm, bottom=2.5cm,
  left=2.8cm, right=2.8cm
}

% ─── Colores institucionales ────────────────────────────────────────────────
\definecolor{uninorteazul}{RGB}{47, 62, 117}
\definecolor{uninorteclaro}{RGB}{220, 230, 250}
\definecolor{verdepaso}{RGB}{39, 174, 96}
\definecolor{verdeclaro}{RGB}{212, 245, 225}
\definecolor{amarilloaviso}{RGB}{255, 193, 7}
\definecolor{amarilloclaro}{RGB}{255, 248, 220}
\definecolor{rojopeligro}{RGB}{192, 57, 43}
\definecolor{rojoclaro}{RGB}{250, 220, 215}
\definecolor{grisclaro}{RGB}{245, 245, 245}
\definecolor{grismedio}{RGB}{180, 180, 180}
\definecolor{codigofondo}{RGB}{30, 30, 30}
\definecolor{codigotexto}{RGB}{220, 220, 170}
\definecolor{marcoimg}{RGB}{100, 120, 180}

% ─── Hyperref ───────────────────────────────────────────────────────────────
\hypersetup{
  colorlinks=true,
  linkcolor=uninorteazul,
  urlcolor=uninorteazul,
  pdftitle={Instructivo - Sistema de Reportes Uninorte},
  pdfauthor={Universidad del Norte}
}

% ─── Encabezado y pie de página ─────────────────────────────────────────────
\pagestyle{fancy}
\fancyhf{}
\fancyhead[L]{\small\color{uninorteazul}\textbf{Sistema de Reportes — Universidad del Norte}}
\fancyhead[R]{\small\color{grismedio}Guía de Instalación}
\fancyfoot[C]{\small\color{grismedio}\thepage}
\renewcommand{\headrulewidth}{0.4pt}
\renewcommand{\footrulewidth}{0pt}

% ─── Títulos de sección ─────────────────────────────────────────────────────
\titleformat{\section}[block]
  {\Large\bfseries\color{uninorteazul}}
  {\colorbox{uninorteazul}{\color{white}\,\thesection\,}}{0.8em}{}
  [\vspace{-0.3em}\color{uninorteazul}\rule{\linewidth}{1.5pt}\vspace{0.3em}]

\titleformat{\subsection}[block]
  {\large\bfseries\color{uninorteazul!80!black}}
  {\thesubsection}{0.6em}{}

% ─── Cajas de entorno ───────────────────────────────────────────────────────

\newtcolorbox{paso}[2][]{
  enhanced, breakable,
  colback=verdeclaro,
  colframe=verdepaso,
  fonttitle=\bfseries\color{white},
  coltitle=white,
  colbacktitle=verdepaso,
  title={\faCheckCircle\quad Paso #2},
  attach boxed title to top left={yshift=-2mm, xshift=4mm},
  boxed title style={colback=verdepaso, rounded corners},
  rounded corners, arc=4pt,
  left=6pt, right=6pt, top=8pt, bottom=6pt,
  #1
}

\newtcolorbox{aviso}[1][Importante]{
  enhanced,
  colback=amarilloclaro,
  colframe=amarilloaviso,
  fonttitle=\bfseries,
  coltitle=black,
  colbacktitle=amarilloaviso,
  title={\faExclamationTriangle\quad #1},
  attach boxed title to top left={yshift=-2mm, xshift=4mm},
  boxed title style={colback=amarilloaviso, rounded corners},
  rounded corners, arc=4pt,
  left=6pt, right=6pt, top=8pt, bottom=6pt,
}

\newtcolorbox{solucion}[1][Si algo falla]{
  enhanced,
  colback=rojoclaro,
  colframe=rojopeligro,
  fonttitle=\bfseries\color{white},
  coltitle=white,
  colbacktitle=rojopeligro,
  title={\faTimesCircle\quad #1},
  attach boxed title to top left={yshift=-2mm, xshift=4mm},
  boxed title style={colback=rojopeligro, rounded corners},
  rounded corners, arc=4pt,
  left=6pt, right=6pt, top=8pt, bottom=6pt,
}

\newtcolorbox{info}[1][Dato útil]{
  enhanced,
  colback=uninorteclaro,
  colframe=uninorteazul,
  fonttitle=\bfseries\color{white},
  coltitle=white,
  colbacktitle=uninorteazul,
  title={\faInfoCircle\quad #1},
  attach boxed title to top left={yshift=-2mm, xshift=4mm},
  boxed title style={colback=uninorteazul, rounded corners},
  rounded corners, arc=4pt,
  left=6pt, right=6pt, top=8pt, bottom=6pt,
}

\newtcolorbox{comando}{
  enhanced,
  colback=codigofondo,
  colframe=grismedio,
  coltext=codigotexto,
  fontupper=\ttfamily\small,
  rounded corners, arc=3pt,
  left=10pt, right=10pt, top=6pt, bottom=6pt,
  boxrule=0.5pt,
}

% ─── Comando para capturas de pantalla enmarcadas ───────────────────────────
\newcommand{\imgoplaceholder}[1]{%
  \begin{center}\color{grismedio}\itshape
    \faImage\;\; Captura de pantalla no disponible en esta copia.\\[2pt]
    \small Consulte esta pantalla directamente en la aplicación (archivo: \texttt{#1}).
  \end{center}%
}

\newcommand{\captura}[2][]{%
  \begin{tcolorbox}[
    enhanced,
    colback=white,
    colframe=marcoimg,
    boxrule=1pt,
    arc=6pt,
    drop shadow={opacity=0.18, shadow xshift=3pt, shadow yshift=-3pt, fill=uninorteazul!30},
    left=4pt, right=4pt, top=4pt, bottom=0pt,
    #1
  ]
    \IfFileExists{#2}{\includegraphics[width=\linewidth]{#2}}{\imgoplaceholder{#2}}%
  \end{tcolorbox}%
}

\newcommand{\capturaConTitulo}[3][]{%
  \begin{tcolorbox}[
    enhanced, breakable,
    colback=white,
    colframe=marcoimg,
    fonttitle=\small\bfseries\color{white},
    colbacktitle=uninorteazul,
    coltitle=white,
    title={\faDesktop\quad #2},
    attach boxed title to top left={yshift=-2mm, xshift=4mm},
    boxed title style={colback=uninorteazul, rounded corners},
    boxrule=1pt,
    arc=6pt,
    drop shadow={opacity=0.18, shadow xshift=3pt, shadow yshift=-3pt, fill=uninorteazul!30},
    left=4pt, right=4pt, top=8pt, bottom=0pt,
    #1
  ]
    \IfFileExists{#3}{\includegraphics[width=\linewidth]{#3}}{\imgoplaceholder{#3}}%
  \end{tcolorbox}%
}

% ─── Documento ──────────────────────────────────────────────────────────────
\begin{document}

% ════════════════════════════════════════════════════════
%  PORTADA
% ════════════════════════════════════════════════════════
\begin{titlepage}
  \pagecolor{uninorteazul}
  \color{white}
  \centering
  \vspace*{2cm}

  {\fontsize{72}{80}\selectfont \faFileMedical}

  \vspace{1cm}
  {\fontsize{11}{13}\selectfont\scshape Universidad del Norte}

  \vspace{0.4cm}
  \rule{10cm}{0.5pt}
  \vspace{0.6cm}

  {\fontsize{28}{34}\selectfont\bfseries
    Sistema Automático\\[0.3em]
    de Reportes de\\[0.3em]
    Satisfacción
  }

  \vspace{0.8cm}
  \rule{10cm}{0.5pt}
  \vspace{0.8cm}

  {\fontsize{16}{20}\selectfont
    \textbf{Instructivo de Instalación y Uso}\\[0.4em]
    \normalsize Para personas sin conocimientos de programación
  }

  \vfill

  \begin{tcolorbox}[
    colback=white!15!uninorteazul,
    colframe=white,
    width=10cm,
    rounded corners, arc=6pt,
    boxrule=1pt
  ]
    \centering\color{white}
    \small
    \faWindows\; \textbf{Requisito:} Windows 10 u 11\\[4pt]
    \faClock\; \textbf{Tiempo estimado:} 15 a 25 minutos\\[4pt]
    \faDesktop\; \textbf{Nivel:} Sin experiencia en programación
  \end{tcolorbox}

  \vspace{1.5cm}
  {\small Versión 2.0 — 2026}

\end{titlepage}
\nopagecolor

% ════════════════════════════════════════════════════════
%  TABLA DE CONTENIDO
% ════════════════════════════════════════════════════════
\newpage
\tableofcontents
\newpage

% ════════════════════════════════════════════════════════
\section{¿Qué hace este programa?}
% ════════════════════════════════════════════════════════

Este programa convierte los resultados de las encuestas de satisfacción en reportes de Excel completos y profesionales, de forma automática.

En lugar de calcular manualmente promedios, porcentajes y crear gráficos, usted simplemente:

\begin{enumerate}[leftmargin=1.5cm, label=\textcolor{verdepaso}{\textbf{\arabic*.}}]
  \item Sube el archivo de Excel con los datos de la encuesta.
  \item Elige la oficina y el proceso correspondiente.
  \item Hace clic en un botón.
  \item Descarga el reporte terminado con tablas, gráficos y filtros interactivos.
\end{enumerate}

\vspace{0.5em}

\capturaConTitulo{Vista general de la plataforma}{Programa 5.png}

\vspace{0.5em}

\begin{info}[¿Qué contiene el reporte generado?]
  El archivo de Excel que se genera incluye automáticamente:
  \begin{itemize}[leftmargin=0.5cm]
    \item Ficha técnica con tamaño de muestra y margen de error.
    \item Tablas de satisfacción por cada pregunta de la encuesta.
    \item Gráficos de barras para cada atributo evaluado.
    \item Índice de Satisfacción del Cliente (ISC) ponderado por correlación.
    \item Matriz de Importancia vs. Satisfacción.
    \item Filtros desplegables para explorar los datos por grupos.
    \item (Opcional) Los atributos agrupados bajo encabezados de categoría.
  \end{itemize}
\end{info}

\begin{info}[Novedades de esta versión (2.0)]
  Además de generar el reporte, ahora puede:
  \begin{itemize}[leftmargin=0.5cm]
    \item \textbf{Agrupar los atributos por categorías} con un encabezado de color (por ejemplo: ``Crédito Uninorte'', ``Becas'', etc.).
    \item \textbf{Segmentar por población}, cuando la encuesta mezcla grupos de distinto tamaño (por ejemplo, pregrado y posgrado).
    \item \textbf{Agregar una oficina o un proceso nuevo} desde la misma pantalla, sin tocar el código.
  \end{itemize}
  Todas estas funciones se explican en la Sección~\ref{sec:avanzadas}.
\end{info}

% ════════════════════════════════════════════════════════
\section{¿Qué necesito antes de empezar?}
% ════════════════════════════════════════════════════════

Antes de instalar el programa, verifique que su computador tenga lo siguiente:

\vspace{0.5em}

\begin{center}
\begin{tabular}{>{\centering\arraybackslash}m{1.5cm} >{\bfseries}m{4cm} m{8.5cm}}
  \toprule
  \textcolor{uninorteazul}{\Large\faWindows} & Sistema Operativo &
    Windows 10 o Windows 11 (el programa \textbf{no funciona en Mac ni Linux}) \\
  \midrule
  \textcolor{uninorteazul}{\Large\faFileExcel} & Microsoft Excel &
    Excel 2010 o superior instalado en el computador \\
  \midrule
  \textcolor{uninorteazul}{\Large\faWifi} & Internet &
    Conexión activa durante la instalación (solo la primera vez) \\
  \midrule
  \textcolor{uninorteazul}{\Large\faHdd} & Espacio en disco &
    Al menos 1 GB de espacio libre disponible \\
  \bottomrule
\end{tabular}
\end{center}

\vspace{0.8em}

\begin{aviso}[Antes de continuar]
  Si no tiene Microsoft Excel instalado, el programa generará el reporte de Excel normalmente, pero \textbf{los filtros interactivos (segmentadores) no se crearán}. Todo lo demás funcionará sin problema.
\end{aviso}

% ════════════════════════════════════════════════════════
\section{Instalación de Python}
% ════════════════════════════════════════════════════════

Python es el motor que hace funcionar este programa. Si nunca ha escuchado este nombre, no se preocupe — es simplemente un programa gratuito que debe instalar una sola vez.

\begin{paso}{1 — Descargar Python}
  \begin{enumerate}[leftmargin=0.5cm, label=\textcolor{uninorteazul}{\arabic*.}]
    \item Abra su navegador de internet (Chrome, Edge, Firefox, etc.)
    \item Vaya a la siguiente dirección web:\\[4pt]
          \texttt{\href{https://www.python.org/downloads/}{https://www.python.org/downloads/}}
    \item Verá un botón amarillo grande que dice \textbf{``Download Python 3.x.x''}. Haga clic en él.
    \item El archivo se descargará automáticamente (pesa aproximadamente 25 MB).
  \end{enumerate}
\end{paso}

\begin{paso}{2 — Instalar Python}
  \begin{enumerate}[leftmargin=0.5cm, label=\textcolor{uninorteazul}{\arabic*.}]
    \item Abra el archivo que descargó (termina en \texttt{.exe}).
    \item \textbf{MUY IMPORTANTE:} Antes de hacer clic en ``Install Now'', marque la casilla que dice:\\[4pt]
          \colorbox{amarilloclaro}{\quad\faCheckSquare\; \textbf{Add Python to PATH}\quad}\\[4pt]
          Esta casilla aparece en la parte inferior de la primera ventana del instalador.
    \item Haga clic en \textbf{``Install Now''}.
    \item Espere a que termine la instalación (1 a 3 minutos).
    \item Haga clic en \textbf{``Close''} cuando aparezca el mensaje de éxito.
  \end{enumerate}
\end{paso}

\begin{aviso}[¿Cómo sé si Python ya estaba instalado?]
  Presione las teclas \textbf{Windows + R} al mismo tiempo, escriba \texttt{cmd} y presione Enter. En la ventana negra que aparece, escriba:
  \begin{comando}
    python --version
  \end{comando}
  Si aparece algo como \texttt{Python 3.x.x}, ya tiene Python instalado y puede saltar al paso 4.
\end{aviso}

% ════════════════════════════════════════════════════════
\section{Descargar los archivos del programa}
% ════════════════════════════════════════════════════════

\begin{paso}{3 — Obtener los archivos del proyecto}

  Si ya tiene una carpeta con los archivos del proyecto (por ejemplo, si se los enviaron por correo o los descargó de un servidor), sáltese este paso.

  Si necesita descargarlos:
  \begin{enumerate}[leftmargin=0.5cm, label=\textcolor{uninorteazul}{\arabic*.}]
    \item Asegúrese de tener el archivo comprimido del proyecto (\texttt{.zip}).
    \item Haga clic derecho sobre el archivo \texttt{.zip}.
    \item Seleccione \textbf{``Extraer todo...''} o \textbf{``Extract All...''}
    \item Elija una ubicación fácil de encontrar, por ejemplo:\\
          \texttt{C:\textbackslash{}Usuarios\textbackslash{}SuNombre\textbackslash{}Desktop\textbackslash{}Automatizacion-Uninorte}
    \item Haga clic en \textbf{``Extraer''}.
  \end{enumerate}

  Al final, debe tener una carpeta que contenga (entre otros) el archivo \texttt{Cargue.py}.
\end{paso}

\begin{info}[¿Cómo encontrar la ruta de la carpeta?]
  Abra la carpeta del proyecto en el Explorador de Archivos, haga clic en la barra de dirección (arriba, donde aparece la ruta de la carpeta) y copie el texto que aparece. Lo necesitará en el siguiente paso.
\end{info}

% ════════════════════════════════════════════════════════
\section{Abrir la ventana de comandos}
% ════════════════════════════════════════════════════════

La \textbf{ventana de comandos} (también llamada ``terminal'' o ``consola'') es una ventana de texto negro donde usted escribe instrucciones para el computador. No tiene que entender cómo funciona; solo debe copiar y pegar los comandos que se indican en este instructivo.

\begin{paso}{4 — Abrir la terminal en la carpeta del proyecto}
  \begin{enumerate}[leftmargin=0.5cm, label=\textcolor{uninorteazul}{\arabic*.}]
    \item Abra el \textbf{Explorador de Archivos} de Windows.
    \item Navegue hasta la carpeta del proyecto (la que contiene \texttt{Cargue.py}).
    \item En la barra de dirección superior (donde aparece la ruta de la carpeta), haga clic una vez para seleccionar todo el texto.
    \item Borre el texto que aparece y escriba exactamente:
    \begin{comando}
      cmd
    \end{comando}
    \item Presione \textbf{Enter}. Se abrirá una ventana negra ya posicionada dentro de su carpeta del proyecto.
  \end{enumerate}
\end{paso}

\begin{aviso}[Alternativa si el método anterior no funciona]
  Presione \textbf{Windows + R}, escriba \texttt{cmd} y presione Enter. Luego escriba el siguiente comando reemplazando la ruta con la de su carpeta:
  \begin{comando}
    cd "C:\textbackslash{}ruta\textbackslash{}de\textbackslash{}su\textbackslash{}carpeta\textbackslash{}del\textbackslash{}proyecto"
  \end{comando}
  Por ejemplo: \texttt{cd "C:\textbackslash{}Users\textbackslash{}Maria\textbackslash{}Desktop\textbackslash{}Automatizacion-Uninorte"}
\end{aviso}

% ════════════════════════════════════════════════════════
\section{Instalación de las librerías necesarias}
% ════════════════════════════════════════════════════════

Las \textbf{librerías} son pequeños programas adicionales que el sistema necesita para funcionar. Solo deben instalarse una vez. Con la terminal abierta en la carpeta del proyecto, siga los pasos a continuación.

\begin{paso}{5 — Instalar todas las librerías de una sola vez}
  La carpeta del proyecto incluye un archivo llamado \texttt{requirements.txt} que contiene la lista completa de librerías necesarias. Para instalarlas todas juntas, en la ventana negra (terminal) escriba o pegue este único comando y presione \textbf{Enter}:
  \begin{comando}
    pip install -r requirements.txt
  \end{comando}
  Verá bastante texto desplazándose en pantalla. Espere hasta que aparezca de nuevo el símbolo \texttt{>} parpadeando. Esto puede tomar entre 3 y 8 minutos la primera vez.
\end{paso}

\begin{info}[¿Qué son estas librerías y para qué sirven?]
  El comando anterior instala automáticamente todas estas piezas:
  \begin{itemize}[leftmargin=0.5cm]
    \item \textbf{streamlit}: Convierte el programa en una página web que se ve en el navegador.
    \item \textbf{pandas} y \textbf{numpy}: Leen y hacen los cálculos con los datos de la encuesta.
    \item \textbf{openpyxl} y \textbf{xlsxwriter}: Leen y crean los archivos de Excel (.xlsx).
    \item \textbf{pywin32}: Permite a Excel crear los filtros interactivos (solo Windows).
    \item \textbf{textwrap3}, \textbf{streamlit-aggrid}, \textbf{matplotlib}, \textbf{seaborn}: Ajustan el texto y muestran tablas y gráficos.
  \end{itemize}
\end{info}

\begin{aviso}[Si el comando anterior no funciona]
  Si aparece un error diciendo que no encuentra \texttt{requirements.txt}, es porque la terminal no está abierta dentro de la carpeta del proyecto. Vuelva al \textbf{Paso 4}. Como alternativa, puede instalar todo escribiendo este comando (una sola línea larga):
  \begin{comando}
    pip install streamlit pandas numpy openpyxl xlsxwriter pywin32 textwrap3 streamlit-aggrid matplotlib seaborn
  \end{comando}
\end{aviso}

\begin{solucion}[Si aparece el error ``pip no se reconoce como comando'']
  Esto significa que Python no se instaló correctamente con la opción PATH. Soluciones:
  \begin{enumerate}[leftmargin=0.5cm, label=\textcolor{rojopeligro}{\arabic*.}]
    \item Desinstale Python desde el Panel de Control de Windows.
    \item Vuelva al \textbf{Paso 2} y asegúrese de marcar la casilla \textbf{``Add Python to PATH''}.
    \item Repita la instalación.
  \end{enumerate}
\end{solucion}

% ════════════════════════════════════════════════════════
\section{Ejecutar el programa}
% ════════════════════════════════════════════════════════

\begin{paso}{6 — Iniciar el sistema}
  Con la terminal abierta \textbf{dentro de la carpeta del proyecto}, escriba o pegue el siguiente comando y presione \textbf{Enter}:
  \begin{comando}
    python -m streamlit run Cargue.py
  \end{comando}
  \vspace{0.3em}
  Verá un mensaje similar a este en la terminal:
  \begin{comando}
    You can now view your Streamlit app in your browser.\\
    \\
    \quad Local URL:  http://localhost:8501\\
    \quad Network URL: http://192.168.x.x:8501
  \end{comando}
  \vspace{0.3em}
  En pocos segundos, \textbf{su navegador se abrirá automáticamente} con la interfaz del programa lista para usar.
\end{paso}

\begin{info}[¿Qué hace ese comando?]
  Le está diciendo al computador: ``Usa Python para ejecutar el programa principal llamado Cargue.py''. La terminal quedará activa mientras el programa esté corriendo — esto es normal. \textbf{No cierre la ventana negra} mientras usa el sistema.
\end{info}

\begin{aviso}[Para detener el programa]
  Cuando termine de usarlo, vaya a la terminal (ventana negra) y presione las teclas \textbf{Ctrl + C} al mismo tiempo. El programa se detendrá.
\end{aviso}

% ════════════════════════════════════════════════════════
\section{Cómo usar el programa paso a paso}
% ════════════════════════════════════════════════════════

Una vez que el programa esté corriendo en el navegador, verá una interfaz con varios pasos. A continuación se describe cada uno con capturas de pantalla reales de la plataforma.

% ──────────────────────────────────────────────────────
\subsection{Paso A — Seleccionar el tipo de procesamiento}
% ──────────────────────────────────────────────────────

Al abrir la aplicación, lo primero que aparece es el título \textbf{``Exportador Automático de Datos a Excel''} y una pregunta sobre cómo están organizados sus datos. Existen tres opciones:

\begin{center}
\begin{tabular}{>{\bfseries\color{uninorteazul}}m{3cm} m{10cm}}
  \toprule
  Opción & Cuándo usarla \\
  \midrule
  Procesar & Sus datos tienen una fila por respuesta, con las calificaciones escritas como texto (por ejemplo, ``Satisfecho'', ``Insatisfecho''). \\
  \addlinespace
  Pivotear & Sus datos tienen una fila por caso y las preguntas están en columnas separadas de Caso, Pregunta y Respuesta. \\
  \addlinespace
  Qualtrics & Sus datos provienen directamente de la plataforma Qualtrics (exportación directa). \\
  \bottomrule
\end{tabular}
\end{center}

\vspace{0.8em}

\capturaConTitulo{Pantalla de inicio — Selección de método y carga del archivo}{Programa 5.png}

% ──────────────────────────────────────────────────────
\subsection{Paso B — Subir el archivo de datos}
% ──────────────────────────────────────────────────────

\begin{enumerate}[leftmargin=1cm, label=\textcolor{verdepaso}{\arabic*.}]
  \item Haga clic en el botón \textbf{``Browse files''} o \textbf{``Examinar archivos''}.
  \item Busque y seleccione su archivo de Excel (\texttt{.xlsx}) con los datos de la encuesta.
  \item El programa procesará el archivo y mostrará una vista previa de los datos en pantalla.
\end{enumerate}

\vspace{0.5em}

\capturaConTitulo{Archivo cargado exitosamente — Vista previa de los datos}{Programa 4.png}

\vspace{0.5em}

Una vez cargado correctamente, verá el mensaje \textbf{``Archivo cargado exitosamente''} en verde y una tabla con los primeros datos de su encuesta.

\begin{aviso}[Formato del archivo de Excel]
  El archivo debe cumplir con estas condiciones mínimas:
  \begin{itemize}[leftmargin=0.5cm]
    \item Estar en formato \textbf{.xlsx} (no .xls, no .csv).
    \item La primera fila debe contener los \textbf{nombres de las columnas} (encabezados).
    \item Las calificaciones de satisfacción deben usar la escala \textbf{1 al 5}.
  \end{itemize}
\end{aviso}

% ──────────────────────────────────────────────────────
\subsection{Paso C — Seleccionar la oficina y configurar parámetros}
% ──────────────────────────────────────────────────────

Justo debajo de los datos cargados aparecerá la sección \textbf{``Seleccionar oficina y parámetros''}:

\begin{enumerate}[leftmargin=1cm, label=\textcolor{verdepaso}{\arabic*.}]
  \item En el menú \textbf{``Selecciona la oficina''}, elija la oficina correspondiente a la encuesta.
  \item Aparecerá un segundo menú con los \textbf{procesos disponibles} para esa oficina. Seleccione el proceso.
  \item Complete los campos del formulario:
\end{enumerate}

\vspace{0.5em}

\begin{center}
\begin{tabular}{>{\bfseries}m{5cm} m{8.5cm}}
  \toprule
  Campo & Qué ingresar \\
  \midrule
  Nombre del archivo de salida & El nombre que quiere darle al Excel generado. Sin extensión \texttt{.xlsx}. \\
  \addlinespace
  Número de población & La cantidad total de personas que podían responder (no los que respondieron, sino el universo completo). \\
  \addlinespace
  Período de la encuesta & El período correspondiente. Ejemplo: \texttt{2025-1} o \texttt{Primer semestre 2025}. \\
  \bottomrule
\end{tabular}
\end{center}

\vspace{0.8em}

\capturaConTitulo{Selección de oficina y configuración de parámetros del reporte}{Porgrama 3.png}

\vspace{0.5em}

\begin{info}[Oficina no encontrada en la lista]
  Si su oficina no aparece en el menú, seleccione \textbf{``Oficina Genérica / Personalizada''} e ingrese manualmente el nombre de su oficina en el campo que aparecerá.
\end{info}

% ──────────────────────────────────────────────────────
\subsection{Paso D — Seleccionar las columnas}
% ──────────────────────────────────────────────────────

Esta sección es la más importante. El programa detecta automáticamente las columnas de preguntas, pero usted debe confirmar y complementar la selección:

\vspace{0.5em}

\capturaConTitulo{Selección de columnas — Preguntas, satisfacción general y filtros}{Programa 2.png}

\vspace{0.8em}

\begin{enumerate}[leftmargin=1cm, label=\textcolor{verdepaso}{\arabic*.}]
  \item \textbf{Columnas de preguntas:} El programa detecta automáticamente las columnas con calificaciones del 1 al 5 y las selecciona. Puede ajustar la selección si es necesario.
  \item \textbf{Columnas de observaciones:} (Opcional) Seleccione las columnas que contienen comentarios o texto libre de los encuestados.
  \item \textbf{Satisfacción general:} Seleccione la columna que contiene la calificación global de satisfacción. En el ejemplo de la imagen se seleccionó \textit{``El nivel de satisfacción general en relación al servicio obtenido es:''}.
  \item \textbf{Columnas para filtros dinámicos:} (Opcional) Seleccione columnas como \textit{Sede}, \textit{Programa}, \textit{Tipo de caso}, etc. El programa creará listas desplegables en el Excel para segmentar los datos por esos grupos (ver Sección~\ref{sec:avanzadas}).
\end{enumerate}

\begin{info}[Agrupar los atributos por categorías]
  Justo debajo de la lista de columnas hay una opción para \textbf{agrupar los atributos bajo encabezados de categoría}. Es opcional y se explica en la Sección~\ref{sec:avanzadas}.
\end{info}

% ──────────────────────────────────────────────────────
\subsection{Paso E — Generar y descargar el reporte}
% ──────────────────────────────────────────────────────

\vspace{0.5em}

\capturaConTitulo{Filtros dinámicos, gráficas y botones de ejecución}{Programa 1.png}

\vspace{0.8em}

\begin{enumerate}[leftmargin=1cm, label=\textcolor{verdepaso}{\arabic*.}]
  \item (Opcional) En \textbf{``Seleccionar gráficas''} puede elegir qué métricas desea visualizar en el reporte.
  \item Haga clic en el botón \textbf{``Ejecutar función excel\_exportar''}.
  \item Espere entre 10 y 60 segundos dependiendo del tamaño de los datos.
  \item Cuando termine, aparecerá el botón \textbf{``Descargar archivo generado''}. Haga clic en él.
  \item El archivo \texttt{.xlsx} se descargará a su carpeta de Descargas.
\end{enumerate}

\begin{info}[Sobre los filtros desplegables]
  Si seleccionó columnas para filtros dinámicos, al abrir el Excel descargado verá, en la hoja \textbf{T+G} (a la derecha del reporte), un panel llamado \textbf{``FILTROS DINÁMICOS''} con listas desplegables. Al elegir un valor en una lista (por ejemplo, un programa académico), todas las tablas, porcentajes y gráficos del reporte se recalculan mostrando solo ese grupo. Para volver a ver todo, elija la opción \textbf{``(Todos)''}.
\end{info}

% ════════════════════════════════════════════════════════
\section{Funciones adicionales (opcionales)}\label{sec:avanzadas}
% ════════════════════════════════════════════════════════

Las siguientes funciones son \textbf{opcionales}. Si solo necesita el reporte estándar, puede ignorarlas por completo. Se activan en la misma pantalla, antes de generar el reporte.

% ──────────────────────────────────────────────────────
\subsection{Filtros desplegables (segmentar los datos)}
% ──────────────────────────────────────────────────────

Sirven para analizar los datos por grupos (por ejemplo: por programa, por sede o por tipo de estudiante) \textbf{sin generar varios reportes}. Un solo archivo permite ver los resultados de cada grupo.

\begin{enumerate}[leftmargin=1cm, label=\textcolor{verdepaso}{\arabic*.}]
  \item En la sección \textbf{``Seleccionar columnas para filtros dinámicos''}, elija una o más columnas (por ejemplo \textit{Programa}, \textit{Sede}).
  \item Genere el reporte normalmente.
  \item Abra el Excel y vaya a la hoja \textbf{T+G}. A la derecha verá el panel \textbf{``FILTROS DINÁMICOS''}.
  \item Elija un valor en cualquier lista desplegable: todo el reporte se actualiza al instante. Elija \textbf{``(Todos)''} para volver a la vista completa.
\end{enumerate}

\begin{info}[No necesita saber de macros]
  Estos filtros son listas desplegables normales de Excel: funcionan con solo abrir el archivo, sin activar macros ni contenido especial.
\end{info}

% ──────────────────────────────────────────────────────
\subsection{Segmentación por población}
% ──────────────────────────────────────────────────────

Úsela cuando su encuesta reúne \textbf{grupos de tamaño muy distinto} (por ejemplo, 3.000 estudiantes de pregrado y 1.000 de posgrado). Así la ficha técnica y la muestra mínima se calculan correctamente para cada grupo.

\begin{enumerate}[leftmargin=1cm, label=\textcolor{verdepaso}{\arabic*.}]
  \item En \textbf{``Segmentación por población''}, elija la \textbf{columna que identifica el grupo} de cada persona (por ejemplo, la columna \textit{``Usted es:''}).
  \item Escriba el \textbf{número total (N)} de personas de cada grupo en los campos que aparecen.
  \item El sistema suma los grupos y, cuando usted filtra por un grupo en el Excel, la población y la muestra se ajustan solas a ese grupo.
\end{enumerate}

\begin{aviso}[El campo ``Número de población'' se ignora]
  Cuando activa la segmentación por población, el sistema usa la suma de los grupos y \textbf{ya no toma en cuenta} el campo ``Número de población'' de más arriba. Esto es normal.
\end{aviso}

% ──────────────────────────────────────────────────────
\subsection{Agrupar los atributos por categorías}
% ──────────────────────────────────────────────────────

Normalmente las tarjetas de cada atributo (pregunta) aparecen una tras otra. Con esta función puede \textbf{agruparlas bajo encabezados de color} con el nombre que usted elija —por ejemplo, ``Canales de atención'', ``Crédito Uninorte'', ``Becas''— tal como en el informe de la OFE.

\begin{enumerate}[leftmargin=1cm, label=\textcolor{verdepaso}{\arabic*.}]
  \item Active la casilla \textbf{``Agrupar las tarjetas de atributos bajo encabezados de categoría''}.
  \item Indique \textbf{cuántas categorías} quiere crear.
  \item A cada categoría póngale un \textbf{nombre} y seleccione los \textbf{atributos} que le corresponden.
  \item Los atributos que no asigne a ninguna categoría se agruparán solos en una categoría final llamada \textbf{``GENERAL''}.
\end{enumerate}

\begin{info}[Los atributos ``se mueven'', no se duplican]
  La lista de arriba (``Atributos sin categoría'') y las categorías comparten los mismos atributos. Cuando selecciona un atributo dentro de una categoría, ese atributo \textbf{desaparece de la lista de arriba} (se movió allí). Si lo quita de la categoría, \textbf{vuelve} a la lista de arriba. De este modo un atributo nunca queda repetido en dos lugares.
\end{info}

% ──────────────────────────────────────────────────────
\subsection{Agregar una oficina o un proceso nuevo}
% ──────────────────────────────────────────────────────

Si su oficina o su proceso no aparecen en las listas, puede crearlos desde la misma pantalla, sin tocar el código. Busque el apartado \textbf{``Agregar o gestionar oficinas y procesos''} (debajo de los menús de oficina y proceso).

\begin{center}
\begin{tabular}{>{\bfseries\color{uninorteazul}}m{4.5cm} m{9cm}}
  \toprule
  Opción & Qué hace \\
  \midrule
  Nueva oficina & Crea una oficina nueva con su primer proceso. Debe elegir un \textbf{``Formato base''} (el diseño del informe que usará). \\
  \addlinespace
  Agregar proceso a oficina existente & Añade otro proceso a una oficina que ya existe (una oficina puede tener varios procesos). \\
  \addlinespace
  Eliminar personalizados & Borra solo lo que usted creó; las oficinas y procesos originales no se tocan. \\
  \bottomrule
\end{tabular}
\end{center}

\vspace{0.5em}

\begin{info}[¿Se crea un programa nuevo para mi oficina?]
  No. Lo que usted agrega se \textbf{guarda de forma permanente} (sobrevive aunque cierre el programa) y su oficina nueva usa el \textbf{mismo diseño de informe} del ``Formato base'' que eligió. El nombre del proceso es solo una etiqueta que aparece en el reporte. Si no está seguro de qué formato base elegir, deje \textbf{``Oficina Genérica / Personalizada''}.
\end{info}

% ════════════════════════════════════════════════════════
\section{Solución a problemas frecuentes}
% ════════════════════════════════════════════════════════

\begin{solucion}[El navegador no se abre automáticamente]
  Abra su navegador manualmente y escriba en la barra de dirección:
  \begin{comando}
    http://localhost:8501
  \end{comando}
  Presione Enter y el programa debería aparecer.
\end{solucion}

\begin{solucion}[Aparece ``ModuleNotFoundError'' en la terminal]
  Significa que falta instalar una librería. El mensaje indicará el nombre de la librería faltante. Por ejemplo, si dice \texttt{No module named 'xlsxwriter'}, ejecute en la terminal:
  \begin{comando}
    pip install xlsxwriter
  \end{comando}
  Reemplace \texttt{xlsxwriter} por el nombre del módulo que indique el error.
\end{solucion}

\begin{solucion}[El programa se cae o se congela al generar el reporte]
  \begin{enumerate}[leftmargin=0.5cm, label=\textcolor{rojopeligro}{\arabic*.}]
    \item Cierre el programa con \textbf{Ctrl + C} en la terminal.
    \item Asegúrese de que Excel no tenga abierto el mismo archivo que está intentando generar.
    \item Vuelva a ejecutar el programa con el comando del Paso 6.
    \item Intente con un archivo de datos más pequeño para descartar problemas de memoria.
  \end{enumerate}
\end{solucion}

\begin{solucion}[Los segmentadores (filtros) no aparecen en el Excel]
  Los segmentadores requieren que Microsoft Excel esté instalado en el computador. Si usa LibreOffice u otro programa para abrir archivos Excel, los segmentadores no funcionarán. Verifique que:
  \begin{enumerate}[leftmargin=0.5cm, label=\textcolor{rojopeligro}{\arabic*.}]
    \item Tiene Microsoft Excel 2010 o superior instalado.
    \item Está abriendo el archivo con Excel (no con otro programa).
  \end{enumerate}
\end{solucion}

\begin{solucion}[Error ``Address already in use'' al iniciar]
  El programa ya está corriendo en otra ventana. Puede acceder directamente desde el navegador en \texttt{http://localhost:8501}, o cierre todas las terminales y vuelva a ejecutar el Paso 6.
\end{solucion}

% ════════════════════════════════════════════════════════
\section{Resumen rápido de comandos}
% ════════════════════════════════════════════════════════

\begin{center}
\begin{tabular}{>{\bfseries}m{5.5cm} m{8cm}}
  \toprule
  \textcolor{uninorteazul}{Acción} & \textcolor{uninorteazul}{Comando en la terminal} \\
  \midrule
  Instalar todas las librerías & \texttt{pip install -r requirements.txt} \\
  \addlinespace
  Iniciar el programa & \texttt{python -m streamlit run Cargue.py} \\
  \addlinespace
  Ver versión de Python & \texttt{python --version} \\
  \addlinespace
  Detener el programa & Presionar \textbf{Ctrl + C} en la terminal \\
  \bottomrule
\end{tabular}
\end{center}

% ════════════════════════════════════════════════════════
\section{Glosario de términos}
% ════════════════════════════════════════════════════════

\begin{description}[leftmargin=2.5cm, style=nextline]
  \item[\textcolor{uninorteazul}{Terminal / Consola}]
    Ventana de fondo negro donde se escriben comandos de texto para controlar el computador. También llamada ``símbolo del sistema'' o ``command prompt''.

  \item[\textcolor{uninorteazul}{Python}]
    Lenguaje de programación gratuito y de código abierto. Es el motor que ejecuta este programa.

  \item[\textcolor{uninorteazul}{pip}]
    Herramienta que se instala junto con Python. Permite instalar librerías adicionales escribiendo un comando simple.

  \item[\textcolor{uninorteazul}{Librería / Módulo}]
    Conjunto de funciones adicionales que se agregan a Python. Son como ``piezas extra'' que amplían las capacidades del lenguaje.

  \item[\textcolor{uninorteazul}{Streamlit}]
    La librería que transforma el código Python en una página web interactiva que se ve en el navegador.

  \item[\textcolor{uninorteazul}{localhost}]
    Dirección especial que significa ``este mismo computador''. Cuando el programa corre en \texttt{localhost:8501}, la página solo es visible en su propio equipo.

  \item[\textcolor{uninorteazul}{NSP}]
    Nivel de Satisfacción Ponderado. Métrica calculada automáticamente que expresa el nivel de satisfacción en una escala de 0 a 100.

  \item[\textcolor{uninorteazul}{ISC}]
    Índice de Satisfacción del Cliente. Indicador global que combina la importancia de cada atributo con su nivel de satisfacción.

  \item[\textcolor{uninorteazul}{Filtro dinámico / Segmentador}]
    Lista desplegable en Excel (hoja T+G) que permite filtrar los datos por grupos. Al elegir un valor, todas las tablas y métricas del reporte se actualizan automáticamente. Con ``(Todos)'' se vuelve a la vista completa.

  \item[\textcolor{uninorteazul}{Atributo}]
    Cada aspecto evaluado en la encuesta (cada pregunta de satisfacción). En el reporte, cada atributo tiene su propia tarjeta con tabla y gráfico.

  \item[\textcolor{uninorteazul}{Categoría}]
    Grupo con nombre bajo el cual se agrupan varias tarjetas de atributos en el reporte (por ejemplo, ``Becas''). Es una función opcional.

  \item[\textcolor{uninorteazul}{Segmentación por población}]
    Opción para indicar el tamaño total (N) de cada grupo cuando la encuesta reúne poblaciones distintas (por ejemplo, pregrado y posgrado), de modo que la muestra se calcule bien para cada uno.

  \item[\textcolor{uninorteazul}{Correlación / NIP}]
    Medida (entre 0 y 1) de qué tan relacionado está un atributo con la satisfacción general. Se usa como ``peso'' (Nivel de Importancia Ponderado) para calcular el ISC: los atributos más relacionados pesan más.
\end{description}

% ════════════════════════════════════════════════════════
%  PIE DE PÁGINA FINAL
% ════════════════════════════════════════════════════════
\vfill
\begin{center}
  \color{grismedio}
  \rule{0.5\linewidth}{0.4pt}\\[6pt]
  \small Universidad del Norte — Dirección de Tecnología Informática y Comunicaciones\\
  Para soporte técnico, contacte al área de Sistemas de Información.
\end{center}

\end{document}
